\chapter{Marco teórico}

Este capítulo describe los conceptos teóricos en los que se apoya el sistema.
Se organiza en tres bloques: el modelo biológico del FST y las conductas que interesa
clasificar; los fundamentos de redes neuronales y aprendizaje profundo para la
clasificación de conductas; y los elementos de arquitectura del sistema web. Las
técnicas específicas de procesamiento de video utilizadas en el módulo de seguimiento
(sustracción de fondo, umbralización, contornos, línea de agua) se describen en el
capítulo~6, dentro del desarrollo del módulo correspondiente.

\section{El modelo de nado forzado como paradigma preclínico}

\subsection{Depresión y modelos animales}

La depresión mayor es uno de los trastornos más frecuentes a nivel mundial. La OMS
estima que afecta a más de 280 millones de personas y es una de las principales causas
de discapacidad \cite{who_depression_2025}. Para estudiar nuevos tratamientos antes de
probarlos en humanos, los investigadores usan modelos animales que reproducen algunos
aspectos del trastorno.

Para que un modelo animal sea considerado válido, debe cumplir con tres criterios
\cite{slattery2012}: tener validez de constructo (los mecanismos biológicos
involucrados deben ser parecidos a los del trastorno humano), validez de cara
(los síntomas observables deben ser similares a los clínicos) y validez
predictiva (el modelo debe responder a los fármacos que funcionan en pacientes).

\subsection{La prueba de nado forzado (FST)}

La prueba de nado forzado fue propuesta por Porsolt et al.\ en 1977 \cite{porsolt1977}.
Consiste en poner a la rata en un cilindro lleno de agua del que no puede salir, y
registrar cómo se comporta. Al principio el animal intenta escapar, pero después de
un rato empieza a quedarse flotando. Esa conducta pasiva se interpreta como
\textit{desesperanza conductual} (en inglés \textit{behavioral despair}), que se
asocia con la pérdida de motivación propia de la depresión \cite{armario2021}.

La prueba funciona bien para evaluar antidepresivos porque la mayoría de los fármacos
que son efectivos en humanos también reducen el tiempo de inmovilidad en ratas
\cite{cryan2005}. Es uno de los experimentos más citados en neurociencia conductual:
desde el 2000 hay más de 8{,}000 artículos que lo reportan \cite{della_valle2025}.

\subsection{Conductas que se observan y cómo se interpretan}

En el FST se pueden distinguir tres tipos de conducta \cite{armario2021,cryan2005}.
Son mutuamente excluyentes, o sea que en cada momento la rata solo está haciendo una:

\begin{itemize}
    \item \textbf{Inmovilidad}: el animal flota haciendo solo lo necesario para no
          hundirse. Es el indicador principal de desesperanza conductual.

    \item \textbf{Nado activo}: movimientos de paladas que lo desplazan horizontalmente.
          Está relacionado con el sistema serotoninérgico; los antidepresivos tipo
          ISRS (como la fluoxetina) aumentan esta conducta.

    \item \textbf{Escalamiento}: movimientos rápidos hacia las paredes del cilindro,
          con las patas delanteras tocando o rompiendo la superficie del agua. Se asocia
          al sistema noradrenérgico; la desipramina, por ejemplo, la incrementa.
\end{itemize}

Distinguir estas tres conductas es importante porque permite saber sobre qué sistema
neurotransmisor está actuando el fármaco evaluado. Un sistema que solo detecta si la
rata se mueve o no pierde esa información.

\subsection{Protocolo del laboratorio colaborador}

En el laboratorio de la ENMyH, las ratas se meten individualmente en cilindros
transparentes con agua. Hay dos sesiones: 20 minutos el día 1 (línea base) y
5 minutos el día 2 (después del tratamiento), con 24 horas de diferencia. Los videos
se graban desde un costado con cámara de celular, capturando los cuatro cilindros al mismo
tiempo. La iluminación y el nivel del agua varían entre experimentos, y con ratas de
pelaje claro el contraste contra el cilindro puede ser bastante bajo. Estas condiciones
representan el caso de uso principal del sistema y determinarán en qué medida los
algoritmos funcionan correctamente; los requisitos mínimos de calidad del video se
definirán durante las pruebas de validación.

\section{Redes neuronales y aprendizaje profundo}

\subsection{Idea general de una red neuronal}

Una red neuronal es un modelo que aprende a partir de ejemplos. Está formada por
capas de ``neuronas'' (en realidad son solo funciones matemáticas simples) que se
conectan entre sí. Cada neurona recibe varios valores, los pondera con unos pesos,
los suma y aplica una función de activación. La función más usada hoy en día es
ReLU, que devuelve el valor si es positivo y cero si es negativo.

Los pesos no se definen a mano: se aprenden durante el entrenamiento. El proceso es
iterativo: se le da al modelo un ejemplo con su respuesta correcta, se ve cuánto se
equivocó, y se ajustan un poquito los pesos para que la próxima vez lo haga mejor.
Eso se hace con el algoritmo de retropropagación y un optimizador como Adam
\cite{kingma2014}, que es una versión mejorada del descenso por gradiente clásico.

\subsection{Redes convolucionales (CNN)}

Las CNN son redes neuronales diseñadas específicamente para trabajar con imágenes
\cite{lecun1998}. La idea principal es que en vez de conectar cada píxel con cada
neurona (lo que sería inmanejable), se usan filtros pequeños que se deslizan por
toda la imagen. El mismo filtro se aplica en cada posición, lo cual reduce mucho la
cantidad de parámetros y hace que la red pueda detectar un patrón sin importar dónde
esté en la imagen.

Las capas del principio aprenden a detectar cosas simples como bordes o cambios de color.
Las capas más profundas van combinando esas características básicas para detectar formas
más complejas, y al final hay capas completamente conectadas que hacen la clasificación.

\subsection{Transfer learning}

Entrenar una red desde cero requiere muchos datos etiquetados, generalmente cientos de
miles de imágenes. Para este proyecto eso no es viable porque el laboratorio tiene un
número limitado de videos.

La solución es el \textit{transfer learning} \cite{pan2010}: se usa una red que ya fue
entrenada en otro dataset grande (en este caso ImageNet, con más de un millón de
imágenes) y se adapta para la tarea nueva. En la práctica se reemplaza la última
capa del modelo (que antes clasificaba 1{,}000 objetos diferentes) por una capa nueva con
solo 3 salidas (una por conducta), y se re-entrena solo esa parte final con los videos
del laboratorio. Las capas anteriores ya saben detectar bordes, texturas y formas
generales, y eso sí se puede aprovechar.

\subsection{ResNet}

ResNet \cite{he2016} es una arquitectura CNN muy usada como punto de partida para transfer
learning. La novedad que introdujo fue agregar conexiones directas que se ``saltan'' una o
más capas (por eso se llaman \textit{skip connections} o conexiones residuales). Eso
facilita mucho el entrenamiento cuando la red es muy profunda, porque el error puede
fluir más fácilmente hacia atrás sin desvanecerse.

Para este proyecto se considera usar ResNet-18 o ResNet-50 como base del clasificador,
dependiendo de los resultados de las pruebas iniciales.

\subsection{Clasificación de las tres conductas}

% [PENDIENTE: definir la arquitectura final del clasificador en Sprint 4]

El clasificador recibe como entrada el recorte de la imagen alrededor de la rata (el
bounding box), normalizado a un tamaño fijo. La capa de salida produce tres valores
de probabilidad (uno por conducta) y se predice la que tenga mayor valor. Para entrenar
el modelo se necesitan cuadros de video etiquetados por el Dr. \dots,
indicando qué conducta realiza cada rata. Esas etiquetas se obtienen con la herramienta
BORIS y sirven también como gold standard para la validación.

\section{Seguimiento de objetos en video}

% [PENDIENTE: desarrollar esta sección con más detalle una vez definida la
%  implementación final del tracker --- ver cap. 6]

\subsection{Trackers de correlación: CSRT y KCF}

Los trackers de correlación aprenden a reconocer cómo se ve el objeto que siguen y
en cada cuadro buscan la región de la imagen que más se le parece. CSRT \cite{lukezic2017}
es más preciso pero tarda más en procesar cada cuadro. KCF \cite{henriques2015} es
más rápido pero menos robusto ante cambios bruscos de apariencia. En el sistema se
usa CSRT como primera opción, con KCF como respaldo si CSRT no está disponible en la
versión de OpenCV instalada.

\subsection{Estrategia de seguimiento en cascada}

% [PENDIENTE: describir con detalle los tres niveles del fallback y agregar diagrama
%  de flujo --- fig. 3.X]

Para que ningún cuadro quede sin detección, el módulo implementa tres niveles de
respaldo. Primero intenta detectar a la rata por segmentación (sustracción de fondo).
Si eso falla o el resultado no es consistente con el cuadro anterior, usa la predicción
del tracker CSRT. Y si ese también falla, congela la última posición conocida,
asignándole una confianza que va disminuyendo con el tiempo hasta un mínimo.

\subsection{Métricas de evaluación}

% [PENDIENTE: agregar ecuaciones de precisión, recall y F1]

Para validar el sistema se compararán sus predicciones contra las anotaciones manuales
del experto. La métrica principal es la kappa de Cohen ($\kappa$), que mide la
concordancia entre dos evaluadores descontando lo que ocurriría por pura casualidad.
Se busca alcanzar $\kappa > 0.80$, que corresponde a concordancia casi perfecta según
la escala de Landis y Koch. También se calcula el error absoluto medio en segundos
para tener una idea del error acumulado por animal.

\section{Arquitectura del sistema web}

% [PENDIENTE: desarrollar esta sección. Temas: API REST con Flask, procesamiento
%  asíncrono con Redis, Docker, PostgreSQL + SQLAlchemy]

\subsection{Arquitectura cliente-servidor y API REST}

% [PENDIENTE: desarrollar]

\subsection{Procesamiento asíncrono con colas de tareas}

% [PENDIENTE: desarrollar. Incluir diagrama del flujo:
%  subida del video -> encolado -> worker -> resultado -> consulta del cliente]

\subsection{Contenedorización con Docker}

% [PENDIENTE: desarrollar]

\subsection{Persistencia de datos con PostgreSQL}

% [PENDIENTE: referenciar al modelo de datos del capítulo 4]
